% GENERATED FILE. Edit canonical lessons, then run npm run generate:books.
% canonical-source-hash: fnv1a64:ce890c153897e279
% canonical-lessons: TA-C162-kata, TA-C162-tanku, TA-C162-curru, TA-C162-nakar, TA-C162-tukku

\chapter{To cross, To stay, To go around, To move, To lift}
\label{ch:ta-to-cross-to-stay-to-go-around-to-move-to-lift}
\hlchaptermodality{162}

\begin{chapteropening}
\textbf{By the end of this chapter:} \emph{I can say to cross, to stay, to go around, to move and to lift.}

Complete the last lesson of chapter 162: I can say to cross, to stay, to go around, to move and to lift.
\end{chapteropening}

\section[\texorpdfstring{\ta{கட} (kaṭa)}{கட (kaṭa)}]{\ta{கட} (kaṭa) — to cross}

\label{lesson:TA-C162-kata}



\begin{quote}
Before the new one: say the Tamil for to grind, then the Tamil for to mix.
\end{quote}

\subsection*{You'll want to know: \ta{கட}}
\textbf{\ta{கட}} — \emph{kaṭa} — \textquotedblleft{}to cross\textquotedblright{}.

\begin{grammarlens}[title={What you've built}]
One more word for this chapter.
\end{grammarlens}

\subsection*{Your turn}
\begin{itemize}
\raggedright
  \item \emph{Say it:} \emph{kaṭa}
  \item \emph{Say it:} \emph{kaṭa}, once more
  \item \emph{Say it:} say \emph{kaṭa}
  \item \emph{Recall:} say the Tamil for to grind, then the Tamil for to mix, then say \emph{kaṭa} again
\end{itemize}

\begin{tcolorbox}[breakable,colback=teal!4,colframe=teal!35!black,title={Before you move on}]
What does \ta{கட} mean? (\textquotedblleft{}To cross\textquotedblright{}.) Say it once more.
\end{tcolorbox}
\section[\texorpdfstring{\ta{தங்கு} (taṅku)}{தங்கு (taṅku)}]{\ta{தங்கு} (taṅku) — to stay (overnight)}

\label{lesson:TA-C162-tanku}



\begin{quote}
Before the new one: say the Tamil for to mix, then the Tamil for to cross.
\end{quote}

\subsection*{You'll want to know: \ta{தங்கு}}
\textbf{\ta{தங்கு}} — \emph{taṅku} — \textquotedblleft{}to stay (overnight)\textquotedblright{}.

\begin{grammarlens}[title={What you've built}]
One more word for this chapter.
\end{grammarlens}

\subsection*{Your turn}
\begin{itemize}
\raggedright
  \item \emph{Say it:} \emph{taṅku}
  \item \emph{Say it:} \emph{taṅku}, once more
  \item \emph{Say it:} say \emph{taṅku}
  \item \emph{Recall:} say the Tamil for to mix, then the Tamil for to cross, then say \emph{taṅku} again
\end{itemize}

\begin{tcolorbox}[breakable,colback=teal!4,colframe=teal!35!black,title={Before you move on}]
What does \ta{தங்கு} mean? (\textquotedblleft{}To stay (overnight)\textquotedblright{}.) Say it once more.
\end{tcolorbox}
\section[\texorpdfstring{\ta{சுற்று} (cuṟṟu)}{சுற்று (cuṟṟu)}]{\ta{சுற்று} (cuṟṟu) — to go around, to wander}

\label{lesson:TA-C162-curru}



\begin{quote}
Before the new one: say the Tamil for to cross, then the Tamil for to stay.
\end{quote}

\subsection*{You'll want to know: \ta{சுற்று}}
\textbf{\ta{சுற்று}} — \emph{cuṟṟu} — \textquotedblleft{}to go around, to wander\textquotedblright{}.

\begin{grammarlens}[title={What you've built}]
One more word for this chapter.
\end{grammarlens}

\subsection*{Your turn}
\begin{itemize}
\raggedright
  \item \emph{Say it:} \emph{cuṟṟu}
  \item \emph{Say it:} \emph{cuṟṟu}, once more
  \item \emph{Say it:} say \emph{cuṟṟu}
  \item \emph{Recall:} say the Tamil for to cross, then the Tamil for to stay, then say \emph{cuṟṟu} again
\end{itemize}

\begin{tcolorbox}[breakable,colback=teal!4,colframe=teal!35!black,title={Before you move on}]
What does \ta{சுற்று} mean? (\textquotedblleft{}To go around, to wander\textquotedblright{}.) Say it once more.
\end{tcolorbox}
\section[\texorpdfstring{\ta{நகர்} (nakar)}{நகர் (nakar)}]{\ta{நகர்} (nakar) — to move, to shift}

\label{lesson:TA-C162-nakar}



\begin{quote}
Before the new one: say the Tamil for to stay, then the Tamil for to go around.
\end{quote}

\subsection*{You'll want to know: \ta{நகர்}}
\textbf{\ta{நகர்}} — \emph{nakar} — \textquotedblleft{}to move, to shift\textquotedblright{}.

\begin{grammarlens}[title={What you've built}]
One more word for this chapter.
\end{grammarlens}

\subsection*{Your turn}
\begin{itemize}
\raggedright
  \item \emph{Say it:} \emph{nakar}
  \item \emph{Say it:} \emph{nakar}, once more
  \item \emph{Say it:} say \emph{nakar}
  \item \emph{Recall:} say the Tamil for to stay, then the Tamil for to go around, then say \emph{nakar} again
\end{itemize}

\begin{tcolorbox}[breakable,colback=teal!4,colframe=teal!35!black,title={Before you move on}]
What does \ta{நகர்} mean? (\textquotedblleft{}To move, to shift\textquotedblright{}.) Say it once more.
\end{tcolorbox}
\section[\texorpdfstring{\ta{தூக்கு} (tūkku)}{தூக்கு (tūkku)}]{\ta{தூக்கு} (tūkku) — to lift}

\label{lesson:TA-C162-tukku}



\begin{quote}
Before the new one: say the Tamil for to go around, then the Tamil for to move.
\end{quote}

\subsection*{You'll want to know: \ta{தூக்கு}}
\textbf{\ta{தூக்கு}} — \emph{tūkku} — \textquotedblleft{}to lift\textquotedblright{}.

\begin{grammarlens}[title={What you've built}]
One more word for this chapter.
\end{grammarlens}

\subsection*{Your turn}
\begin{itemize}
\raggedright
  \item \emph{Say it:} \emph{tūkku}
  \item \emph{Say it:} \emph{tūkku}, once more
  \item \emph{Say it:} say \emph{tūkku}
  \item \emph{Recall:} say the Tamil for to go around, then the Tamil for to move, then say \emph{tūkku} again
\end{itemize}

\begin{tcolorbox}[breakable,colback=teal!4,colframe=teal!35!black,title={Before you move on}]
What does \ta{தூக்கு} mean? (\textquotedblleft{}To lift\textquotedblright{}.) Say it once more.
\end{tcolorbox}
